\def\ChapterCode{WIS}
\chapter
[\CO{[WIS]} Wissenschaft und Forschung]
{Wissenschaft und Forschung}
\label{chp:WIS}



\ChapterIntro
{%
\SectionUnderDevelopment
}
{%
\begin{description}
  \item[Maintainer:] Sebastian Gabriel
  \item[Autoren:] Diverse
  \item[Reviewer:] Standard-Reviewprozess
  \item[Version:] {\DocVersion} ({\DocVersionInfo})
  \item[SHA1:] {\DocGitCommit}
\end{description}

{\TxTeilkapitelAASS}
}

\section{Von der Studie zur Lehrmeinung}

\subsection{Grundlagen -- Was ist Wissenschaft?}

\subsubsection{Wissensgewinn: Am Anfang war die Frage -- am
Ende auch}

\subsubsection{Welt in Zahlen ausgedrückt: Die Rolle der
Statistik}

\subsubsection{Was sind Studien?}

\subsubsection{Fachgesellschaften}

In der Welt der Wissenschaft enstehen und entstanden sog.
Fachgesellschaften, die sich bestimmten Fachgebieten,
Erkrankungen und Fragestellungen widmen.

\cite{Bahr200701}: \bibentry{Bahr200701}

\TextSlide{Medizinische Fachgesellschaften}
{%
\label{topic:fachgesellschaft}
\index{Fachgesellschaft}
Die Tätigkeit von medizinische Fachgesellschaften umfasst
oft die Forschung, Fortbildung und die Herausgabe von
\EE{Empfehlungen} und Leitlinien zur Diagnostik, Therapie
und Vorbeugung u.\,ä. Diese Empfehlungen sind nicht
\enquote{gesetzlich bindend}, aber geben i.\,d.\,R. den
\EE{Stand der Wissenschaft} wieder. Ausserdem müssen diese,
oft allgemein formulierten, Empfehlungen an die örtlichen
Gegebenheiten angepasst werden. Verschiedene
Fachgesellschaften können -- auch zu den gleichen Themen --
unterscheidliche Meinungen haben und unterschiedliche
Empfehlungen aussprechen.
}{%
\begin{itemize}
   \item Forschung, Fortbildung 
   \item Herausgabe von Empfehlungen für Diagnostik, Therapie und
   Vorbeugung (u.\,a.)
   \item Nicht bindend, aber \enquote{Stand der Wissenschaft}
\end{itemize}
}{}{}{}




\section{Statistik}


\Captionof{table}{Statistische Tests mit lustigen Namen -- Eine Übersicht.}

\ChapterEnd